We conducted a reproducible, protocol-matched comparison of compact CNN, frequency-aware, and state-space detectors for AI-generated still-image detection under a locked no-pretraining recipe.
Using frozen splits and clearly separated metrics---especially \textbf{UFD Macro} versus \textbf{OOD Pooled}---we find that \textbf{LiteSSM-A} offers the most favorable generalization--efficiency operating point among Panel~A models, combining strong ID performance, the highest UFD Macro AUC in that panel, and improved bedroom-domain behavior at moderate model size, with a locked batch-32 throughput of 226 images/s.
It is selected for that trade-off rather than for lowest absolute batch-1 latency.
LiteSSM-B remains a competitive SSM alternative under some aggregates (including a higher OOD Pooled AUC of 0.722).

Equally important are the boundaries that remain: DALL·E detection stays near chance for all Panel~A architectures while Panel~B pretrained references remain strong on the same subset; CelebA-HQ saturation masks bedroom difficulty in pooled ID scores; JPEG Q70 leaves all evaluated compact models nearly unchanged without crowning a frequency-specific winner; and frequency branches do not universally improve SSM or CNN transfer.
Relative to LAID and LOTA~\cite{chivaran2025laid,wang2025lota}, the reusable knowledge is the controlled evidence organization---generator-macro OOD, domain decomposition, worst-generator failure, and measured latency/throughput under one disclosed protocol---not a claim of a new universal lightweight detector.
These results support selecting LiteSSM-A as a practical compact operating point while rejecting narratives of universal, domain-invariant, or broadly compression-robust detection beyond the evaluated settings.

Future work will pre-register a modern-generator augmentation protocol with an independent development split, targeting improved FLUX/SDXL transfer without degrading the frozen UFD Macro baseline, and without treating held-out modern-generator gains as method-selection evidence.
Additional directions include extending multi-seed reporting beyond the LiteSSM-A / EfficientNet-B0 appendix probe and hardening worst-generator monitoring while preserving the compact, fully disclosed protocol used here.
